\documentclass{article}

\usepackage[spanish]{babel}
\usepackage[utf8]{inputenc}
\usepackage[a4paper,headheight=42pt,hoffset=0.6cm, left=1cm, right=2cm, top=4cm, bottom=2cm]{geometry}
\usepackage{graphicx}
\usepackage{fancyhdr}
\usepackage{float}
\usepackage{listings}
\usepackage{parcolumns}
\usepackage{listings}
\usepackage{wrapfig}
\usepackage{lscape}
\usepackage{rotating}
\usepackage{epstopdf}
\usepackage{xcolor}
\definecolor{commentgreen}{RGB}{2,112,10}
\definecolor{eminence}{RGB}{108,48,130}
\definecolor{weborange}{RGB}{255,165,0}
\definecolor{frenchplum}{RGB}{129,20,83}


\usepackage[hang,bf]{caption}

%-----------------------INFORMACIÓN DE MATERIA--------------------------------%
\newcommand{\codMateria}{96.08}
\newcommand{\nombreMateria}{Taller de programación 1}
\newcommand{\cuatrimestre}{1$^{er}$ Cuatrimestre de 2018}
\newcommand{\titulo}{Manual del proyecto}
\newcommand{\grupo}{Grupo 2}
%-----------------------INFORMACIÓN DE MATERIA--------------------------------%

%-----------------------INFORMACIÓN DE ALUMNOS--------------------------------%
\newcommand{\nombreUno}{Bourbon Navarro, Rodrigo}
\newcommand{\padronUno}{96961}
\newcommand{\emailUno}{fedemgp@gmail.com}

\newcommand{\nombreDos}{Gómez Peter, Federico}
\newcommand{\padronDos}{96091}
\newcommand{\emailDos}{fedemgp@gmail.com}

\newcommand{\nombreTres}{Lidjens, Axel}
\newcommand{\padronTres}{95772}
\newcommand{\emailTres}{fedemgp@gmail.com}
%-----------------------INFORMACIÓN DE ALUMNOS--------------------------------%
% Esta función se la llama \unidad{mW} e imprime la unidad dada en imprenta
% con un medio espacio.
\newcommand{\unidad}[1]{\mbox{\,#1}}

%Cuando quieran agregar una imagen no tienen que poner el directorio de la
%imagen, eso se está haciendo acá, simplemente pongan el nombre de la imagen

\newcommand{\imgdir}{img}
\graphicspath{{\imgdir/}}

\newcommand{\addFigure}[4]{
	\begin{figure}
		\centering
		\includegraphics[scale=#1]{#2}
		\caption{#3}
		\label{#4}
	\end{figure}
}

\renewcommand{\lstlistingname}{Código}

\begin{document}
\pagestyle{fancy}
\renewcommand{\sectionmark}[1]{\markboth{}{\thesection\ \ #1}}
\lhead{\includegraphics[width= 4 cm]{logofiuba}}
\chead{}
\rhead{\codMateria - \nombreMateria}
\lfoot{}
\cfoot{\thepage}
\rfoot{}

\newpage
\pagestyle{fancy}
\renewcommand{\sectionmark}[1]{\markboth{}{\thesection\ \ #1}}
\lhead{\includegraphics[width= 2.5 cm]{logofiuba}}
\chead{}
\rhead{\rightmark}
\lfoot{}
\cfoot{}
\rfoot{\thepage}

\begin{titlepage}
\begin{center}
\LARGE{Universidad de Buenos Aires}\\
\LARGE{Facultad de Ingeniería}\\
\vspace{10mm}
\includegraphics[scale=0.3]{fiuba}\\
\vspace{10mm}
\LARGE{\codMateria - \nombreMateria}\\[0.5cm] 
\Large{\cuatrimestre}\\[4cm] 
\Large{Trabajo Práctico Final}\\[0.5cm]
\Large{\titulo}\\[4cm]
\large{\grupo} \\
\large{\nombreUno - \#\padronUno} \\
\large{\nombreDos - \#\padronDos} \\
\large{\nombreTres - \#\padronTres}\\

\end{center}
\end{titlepage}

\tableofcontents
\newpage

\section{Integrantes y división de tareas}
\label{sec:integrantes}	
%

\textbf{Axel Lidjens}

\begin{itemize}
	\item Editor.
	\item Carga del nivel de juego mediante los datos provistos por el editor.
	\item Protocolo de comunicación durante el juego.
	\item Encapsulamiento de animaciones, texturas y texto en clases para un fácil manejo de los mismos.
	\item Manejo de la cámara.
	\item Lógica de contacto de los \textit{worms} para evitar que se muevan entre ellos.
	\item Lógica de respuesta de los \textit{worms} frente a vigas inclinadas.
\end{itemize}

\textbf{Federico Gómez Peter}

\begin{itemize}
	\item Lógica de estados de armas y generación de proyectiles, con su posterior destrucción.
	\item Animación de las armas, mira y barra de carga.
	\item Lógia de estados de \textit{worms}.
	\item Parámetros configurables mediante YAML.
	\item Serialización de estado del juego mediante YAML.
	\item Cierre ordenado del servidor.
\end{itemize}

\textbf{Rodrigo Bourbon Navarro}

\begin{itemize}
	\item Lógica de estados de juego.
	\item Lógia de estados de \textit{worms}.
	\item Generación de equipos.
	\item Interfaz gráfica del cliente y protocolo de comunicación con el servidor hasta el comienzo de las partidas.
	\item Lógica de partidas en simultáneo en el servidor.
	\item Cierre ordenado del servidor.
\end{itemize}

\newpage

\section{Inconvenientes}
\label{sec:inconvenientes}	
\subsection{Serialización de mensajes de forma portable}

Se realizó mediante YAML, pero la performance se vio severamente perjudicada, pasando de un consumo de CPU de 4\% de los clientes a un 50\%, el servidor de 2\% a 25\%.


\subsection{Herramientas}
%\label{sec:herramientas}
%\input{herramientas.tex}
\subsection{Editores de interfaz gráfica}

\begin{itemize}
	\item CLion: IDE de \textit{JetBrains} que permite compilar y debuguear con y sin \textit{valgrind} en base a un \textit{CMake}.
	\item VSCode: editor de texto de \textit{Microsoft} que permite compilar y debuguear, pero sus capacidades son más limitadas que las del CLion.
\end{itemize}

\subsection{Control de versiones}

Se trabajó con un control de versiones \textit{git}, en un repositorio alojado en \textit{GitHub}.

%\section{Configuración}
%\label{sec:config}
%El juego es completamente configurable. Para esto, se requiere modificar el archivo de formato \emph{YAML}, que se encuentra en \texttt{/etc/Worms/serverConfig.yaml}. En este archivo se pueden modificar cuestiones físicas del juego ( velocidad de saltos, de caminar, fuerza del viento) como cuestiones del juego (tiempo de turnos, tiempo para mover el worm después de haber disparado, tiempo de espera hasta el siguiente turno, daño de las distintas armas, si las balas son afectadas por el viento, etc).

Para configurar otros aspectos, como lo puede ser el mapa, las municiones de cada arma o la vida de los gusanos, se puede utilizar el editor.

%
%\section{Forma de uso}
%\label{sec:uso}
%\subsection{Servidor}

Para levantar un servidor, se requiere que se ejecute la siguiente linea:

\lstset{
	language=Bash, 
    commentstyle=\color{commentgreen},
    keywordstyle=\color{eminence},
    stringstyle=\color{red},
    basicstyle=\small\ttfamily, % basic font setting
    emph={int,char,double,float,unsigned,void,bool,NULL},
    emphstyle={\color{blue}},
    escapechar=\&,
    breaklines=true,
    numbers=left,
    numberstyle=\tiny, 
    stepnumber=1, 
    numbersep=-2pt
    % keyword highlighting
    }
\begin{lstlisting}[frame=single, caption=Instalación de dependencias, label=code:openServer]
	worms_server <PORT>
\end{lstlisting}

Donde el puerto debe ser mayor a 1000 si es que no se quiere ejecutar el comando como \texttt{sudo}. Además, si se quiere realizar cambios en la configuración del juego, es necesario hacerlo desde el servidor. 

El archivo \texttt{serverConfig.yaml}, que se encuentra en \texttt{/etc/Worms}, tiene la configuración del juego. En este archivo se pueden modificar parámetros como los tiempos de los turnos (\emph{turnTime}, \emph{extraTurnTime}, \emph{waitForNextTurnTime}, \emph{teleportTime}), como también la distancia de caida que puede realizar el gusano sin recibir daño (\emph{safeFallDistance}), su daño maximo, etc. También se pueden modificar parámetros de las armas, como lo pueden ser el daño de esta, el radio del impacto, el factor de amortiguamiento de las explosiones (cuanto impulso se le asigna al gusano impactado), el tiempo inicial para que explote, la friccion y rebote de la bala ( efecto que se apreciará si es un arma que no explota al primer contacto), etc.

\end{document}